\section{环路群的范畴表示}

\label{s:catloop}

\subsection{概览}

\ssszh 到此为止，我们讨论了约化群的 Beilinson--Bernstein 局部化与范畴表示论。

在本综述的余下部分，我们想转向它们的仿射类比，即从约化群过渡到环路群。这方面的故事要丰富得多，正是局部几何朗兰兹对应的主题。

本节将对环路群及其范畴表示作若干介绍性的评论，为第 \ref{s:affbbloc} 节讨论朗兰兹对应作准备。

\subsection{环路群}

\ssszh 与之前一样，仍用 $G$ 记代数闭域 $k$（特征为零）上的连通约化群。

本节的主角不是 $G$，而是它的环路群 $LG$。初步地说，正如 $G$ 由元素在 $k$ 中的某些矩阵构成（取一个忠实表示即可），$LG$ 由同样形状的矩阵构成，只是元素换成了 Laurent 级数域 $k(\!(t)\!)$。更仔细地说，由于我们需要把 $LG$ 理解成代数几何的对象，它是某个群归纳概形，其闭点正如上面所述。

\begin{ex}\label{ex:ga} 当 $G$ 是（非约化的）加法群 $\mathbb{G}_a$ 时，有一个显式的构造，它涵盖了大部分有代表性的想法。此时 $\mathbb{G}_a(k) \simeq k$，即一维仿射空间。如上所述，我们希望 $L\mathbb{G}_a(k) \simeq k(\!(t)\!)$。基本想法是：$k(\!(t)\!)$ 仍是无限维仿射空间；事实上，把 Laurent 级数极点的阶数作界，它就成为一个射有限维仿射空间的递增并。显式地，作为向量空间
%
\[
  k(\!(t)\!) = \varinjlim_n \varprojlim_m t^{-n} \cdot k[[t]] / t^{-n+m}
  \cdot k[[t]],
\]
%
而描述其归纳概形结构用的是同一个公式。也就是说，固定 $n$，得到一个射有限维仿射空间；显式地说，它是可数多个变量的多项式环的谱，而这些变量正是级数 Taylor 系数的函数。让 $n$ 变动，就得到无限维仿射空间沿闭嵌入的递增并，即一个归纳概形。
\end{ex}

\begin{re} 一般地，$LG$ 可以用它的点函子简洁地描述，定义如下。对任意簇 $X$，考虑从仿射簇到集合的函子 $LX$，它在任意 $k$-代数 $R$ 上的 $R$-点取为
	 %
	\[
	     LX(R) = X(R(\!(t)\!) ).
	\]
	%
	这个函子对一般的簇未必表现良好；但若 $X$ 是仿射的，它就由某个归纳无限型的归纳概形表示。取 $X = G$，就得到环路群 $LG$。
\end{re}

给出了形式定义之后，我们仿照例 \ref{ex:ga} 的精神，试着给出更可动手的描述。回忆前几节讨论过的 $G(\Fq)$ 与 $G(k)$ 之间的类比，更一般地，有限 $G(\Fq)$-集与有限维 $G(k)$-簇之间的类比。当下情形的一条基本指导原则，是把 $p$-adic 群 $G(\Fq(\!(t)\!))$ 与环路群 $LG$ 作类比。

\begin{re} 有些作者用 $p$-adic 群一词专指 $G$ 在 $\mathbb{Q}_p$ 或其有限扩张中的点；这里我们略微宽松地用这个词指任意非 Archimedes 局部域（例如 $\Fq(\!(t)\!)$）中的点。 \end{re}


回忆：$p$-adic 群带上其自然拓扑是归纳射有限集。显式地说，取定任一紧开子群 $K$，即满足 $K \cap G(\Fq[[t]])$ 在 $K$ 与 $G(\Fq[[t]])$ 中都是有限指数的子群。那么 $K$ 自然地是有限群的逆向极限，而整个群分解为它的 $K$-陪集的不交并。具体到 $G(\Fq[[t]])$，相应的逆向极限来自幂级数的拓扑，即
%
\begin{equation}
\label{e:invlim}
      G(\Fq[[t]]) = \varprojlim_n G(\Fq[t]/(t^n)).
\end{equation}
%


类似地，可以理解 $LG$ 的紧开子群。具体地说，有弧群 $L^+G$，它是无限型的群概形，$k$-点为 $G(k[[t]])$。它的定义与 \eqref{e:invlim} 类似，是有限维代数群的逆向极限。$LG$ 的子群概形 $K$ 称作紧开的，如果 $K \cap L^+G$ 在 $K$ 与 $L^+G$ 中都是有限余维的。此时 $K$ 自然是代数群的逆向极限，而陪集空间 $LG/K$ 是归纳有限型的归纳概形。我们要强调：在有限集与有限维簇的类比之下，这对应于 $p$-adic 群情形中陪集的可数离散并。



\subsection{环路群上的 D-模与范畴作用}

\ssszh 接下来想讨论 $LG$ 上 D-模的卷积代数。由于这是无限维的对象，需要多加小心。

它的定义可以这样得到。回忆我们把 $\Dmod(G)$ 的模看成类似于 $G(\Fq)$ 的表示；特别地，D-模的卷积代数本身就类比于群代数 $\Fun(G(\Fq))$。

我们同样希望 $\Dmod(LG)$ 的模表现得像 $p$-adic 群 $G(\Fqt)$ 的表示。后者是拓扑群，相应地就要用光滑表示。\footnote{显式地说，复向量空间 $V$ 上的一个表示是光滑的，如果作用映射 $G(\Fqt) \times V \rightarrow V$ 关于 $V$ 上的离散拓扑连续，即每个向量的稳定化子都是开子群。对不熟悉 $p$-adic 表示论的读者，这里的要点是：这个条件对应李群表示中的类似条件，那里要求作用在流形的意义下适当光滑。不过在 $p$-adic 情形中，群上的拓扑与向量空间上的拓扑（不严格地说）是「互素的」。实际上，研究自守形式时自然出现的正是光滑表示。}因此，$G(\Fqt)$ 的模并不等同于其群代数的模，而应当是其 Hecke 代数
$$\on{Dist}(G(\Fqt)) $$
即紧支分布代数的模。具体地说，它可以写成余极限
%
\begin{equation} \label{e:padicgroups}
\on{Dist}(G(\Fqt))\simeq \varinjlim_K  \on{Dist}( K \backslash G(\Fqt) / K),
\end{equation}
%
其中 $K$ 取遍紧开子群，而每个 $\on{Dist}( K \backslash G(\Fqt) / K)$ 不过是双陪集上 delta 分布的有限线性组合，即把函数沿某个双陪集 $KgK$（关于某个 Haar 测度）积分的分布。


有了这些准备，环路群上的 D-模范畴定义为
%
\[
\Dmod(LG) \simeq \varinjlim \Dmod(K \backslash LG / K),
\]
%
其中 $K$ 取遍紧开子群，余极限在 $\DGCat$ 中沿 $*$-拉回函子取。剩下来要描述的是，对固定的 $K$，范畴 $\Dmod(K\backslash LG/K)$ 是什么。回忆 $LG/K$ 是归纳有限型的归纳概形，故还可以把它写成递增并
%
\[
   LG/K = \bigcup_i X_i,
\]
%
其中每个 $X_i$ 是有限型的闭子概形，并且是 $K$-不变的。由于 $K$ 通过一个有限维商 $Q_i$（其核为射幂幺群）作用在每个 $X_i$ 上，有
%
\[
   \Dmod(X_i)^K \simeq \Dmod(X_i)^{Q_i} \simeq \Dmod(Q_i \backslash X_i),
\]
%
也就是说，我们面对的是通常的有限维叠上的 D-模。于是令
%
\[
  \Dmod(K \backslash LG/K) \simeq \varinjlim \Dmod( Q_i \backslash X_i).
\]
%

这里的要点是：粗略地说，$\Dmod(LG)$ 中的对象类似于 $G(\Fqt)$ 上的紧支分布。更精确地说，由例 \ref{e:colimitcats}，$\Dmod(LG)$ 是紧生成的，其紧对象由某个 $K$ 的 $LG/K$ 的子簇上凝聚 D-模有界复形的 $*$-扩张给出。


这个构造的确切细节对我们并不那么重要，不过下面给出两条帮助定位的注。

\begin{re}第一，可以把上述余极限 \eqref{e:padicgroups} 中所有可能的 $K$ 换成任一共尾子集。一个具体的取法是 $L^+G$ 的同余子群 $K_n$（$n \geqslant 1$），我们只需在 $k$-点的层面用正合列来描述它
	%
	\[
	      1 \rightarrow K_n \rightarrow L^+G = G(k[[t]]) \rightarrow
	      G(k[[t]]/(t^n)) \rightarrow 1.
	\]
	%

注意 $K_1$ 是 $L^+G$ 的射幂幺根，由此可知任意充分小的紧开子群 $K$ 都是射幂幺的。特别地，对两个这样的射幂幺子群 $K' \subset K$，拉回
%
\[
    \Dmod(K \backslash LG / K) \rightarrow \Dmod(K' \backslash LG / K')
\]
%
是全忠实的嵌入，参见注 \ref{r:invariantstwice}。
\end{re}



\begin{re}上述构造看起来或许有些凑合，但应当指出：它是为更大一类无限维簇（称作平缓归纳概形）给出 D-模定义的做法中的一个特例 \cite{rdm}。

事实上，这一无限维设定中有一处基本的微妙之处：D-模范畴有两个对偶的版本。一个带有 $*$-推前的函子性，另一个带有 $!$-拉回的函子性。上面的表述对应于带有 $*$-推前的那个。这正是 $G(\Fqt)$ 上函数与分布之区别的类比。与此相关，把维数主丛（dimension torsor）平凡化（即固定一个紧开子群）之后，$LG$ 的这两个对偶范畴可以等同起来。
\end{re}



\ssszh 与有限维情形一样，$\Dmod(LG)$ 在卷积下是幺半 dg 范畴。特别地，如前可以定义它的模范畴
%
\[
      \Dmod(LG)\mod.
\]
%

下面是一些表示的基本例子。

\begin{ex} 若 $X$ 是被 $LG$ 作用的平缓归纳概形，其 D-模范畴 $\Dmod(X)$ 通过卷积带有 $\Dmod(LG)$ 的典范作用。
\end{ex}


\begin{ex}\label{ex:kmmods} 给定环路群的中心扩张
	\[
	1 \rightarrow \mathbb{G}_m \rightarrow \widetilde{LG} \rightarrow LG
	\rightarrow 1,
	\]
	%
	过渡到李代数，就得到环路代数的中心扩张，即仿射李代数
	\[
	 0 \rightarrow k \cdot 1 \rightarrow \gk \rightarrow L \fg \rightarrow
	 0.
	\]
	%
	光滑模的导出范畴
	$
	    \gk\mod
	$
	带有 $\Dmod(LG)$ 的典范作用；这来自 $\widetilde{LG}$ 在 $\widehat{\fg}$ 上的伴随作用通过 $LG$ 分解这一事实。按定义，相应阿贝尔范畴的对象是 $\gk$ 的表示，其中 $1$ 作用为恒等，且每个向量都被某个紧开子代数零化。依照 Frenkel--Gaitsgory \cite{fglocaff}，完整的 dg 范畴是对朴素无界导出范畴作小幅修正，使得 Verma 模、更一般地紧开子代数有限维表示的诱导，都是紧生成元。

	这些细节对我们并不重要，只需指出：这一修正的存在与用处，源于该阿贝尔范畴具有无限的上同调维数。与此相关，重整化后的导出范畴带有 $t$-结构，并在有下界部分与通常的导出范畴典范地等同。也就是说，重整化只在同调次数 $-\infty$ 处带来差别。类似地，范畴 $\Dmod(LG)$ 也带有 $t$-结构，它在整体平移的意义下是典范的，\footnote{这对应 $G(\Fqt)$ 上 Haar 测度在相差一个常数因子意义下的唯一性。}并具有这样的性质：除同调次数 $-\infty$ 处的问题外，$\Dmod(LG)$ 与其心的朴素无界导出范畴一致。
\end{ex}
